\renewcommand{\thesection}{\arabic{section}}\renewcommand{\theHsection}{chapter.\arabic{section}}\setcounter{section}{6}\section{积分学的应用}
\noindent\begin{minipage}{\textwidth}\ProblemHeading{例题}{7.1}
求抛物线 $y^2=2x$ 与直线 $x-y=4$ 围成的区域的面积。
\end{minipage}\par\vfill
\noindent\begin{minipage}{\textwidth}\ProblemHeading{例题}{7.2}
求极坐标方程 $r=a(1+\cos\theta)$（$a>0$）表示的心脏线所围区域的面积。
\end{minipage}\par\vfill
\noindent\begin{minipage}{\textwidth}\ProblemHeading{例题}{7.3}
求星形线 $x=a\cos^3t$、$y=a\sin^3t$（$0\le t\le2\pi$，$a>0$）的弧长。
\end{minipage}\par\vfill
\newpage
\noindent\begin{minipage}{\textwidth}\ProblemHeading{例题}{7.4}
求椭圆 $x^2/a^2+y^2/b^2=1$（$a>b>0$）的周长，可用定积分表示。
\end{minipage}\par\vfill
\noindent\begin{minipage}{\textwidth}\ProblemHeading{例题}{7.5}
求极坐标心脏线 $r=a(1+\cos\theta)$（$a>0$）的弧长。
\end{minipage}\par\vfill
\noindent\begin{minipage}{\textwidth}\ProblemHeading{例题}{7.6}
求两个圆柱体 $x^2+y^2\le a^2$、$x^2+z^2\le a^2$（$a>0$）相交部分的体积。
\end{minipage}\par\vfill
\newpage
\noindent\begin{minipage}{\textwidth}\ProblemHeading{例题}{7.7}
用积分求半径为 $a>0$ 的球体的体积。
\end{minipage}\par\vfill
\noindent\begin{minipage}{\textwidth}\ProblemHeading{例题}{7.8}
用积分计算半径为 $a>0$ 的球面的面积。
\end{minipage}\par\vfill
\noindent\begin{minipage}{\textwidth}\ProblemHeading{例题}{7.9}
求椭圆 $x^2/a^2+y^2/b^2=1$（$a>b>0$）绕 $x$ 轴旋转得到的椭球面的面积。
\end{minipage}\par\vfill
\newpage
\noindent\begin{minipage}{\textwidth}\ProblemHeading{习题}{7.10}
求下列曲线围成的区域的面积：
\SubQ{1}{$ax=y^2$、$ay=x^2$，$a>0$。}
\SubQ{2}{$y=2x-x^2$、$x+y=2$。}
\SubQ{3}{$y=x$、$y=x+\sin^2x$，$0\le x\le\pi$。}
\end{minipage}\par\vfill
\noindent\begin{minipage}{\textwidth}\ProblemHeading{习题}{7.11}
求下列曲线围成的区域的面积：
\SubQ{1}{$ax^2+2bxy+cy^2=1$，$ac-b^2>0$、$c>0$。}
\SubQ{2}{三叶线 $r=a\sin(3\theta)$，$a>0$。}
\SubQ{3}{双纽线 $r^2=a^2\cos(2\theta)$，$a>0$。}
\end{minipage}\par\vfill
\noindent\begin{minipage}{\textwidth}\ProblemHeading{习题}{7.12}
设出现的半轴长 $a,b,c$ 均为正数。
\SubQ{1}{求椭圆域 $x^2/a^2+y^2/b^2\le1$ 的面积。}
\SubQ{2}{求椭球体 $x^2/a^2+y^2/b^2+z^2/c^2\le1$ 的体积。}
\end{minipage}\par\vfill
\newpage
\noindent\begin{minipage}{\textwidth}\ProblemHeading{习题}{7.13}
计算下列曲线的弧长：
\SubQ{1}{$x=e^t\cos t$，$y=e^t\sin t$，$0\le t\le2\pi$。}
\SubQ{2}{$x=a(\cos t+t\sin t)$，$y=a(\sin t-t\cos t)$，$0\le t\le2\pi$、$a>0$。}
\SubQ{3}{$x=(c^2/a)\cos^3t$，$y=(c^2/b)\sin^3t$，$0\le t\le2\pi$，$a>b>0$、$c^2=a^2-b^2$。}
\end{minipage}\par\vfill
\noindent\begin{minipage}{\textwidth}\ProblemHeading{习题}{7.14}
计算下列曲线的弧长：
\SubQ{1}{$x=\sin t$，$y=t$，$z=1-\cos t$，$0\le t\le2\pi$。}
\SubQ{2}{$x=t$，$y=3t^2$，$z=6t^3$，$0\le t\le2$。}
\SubQ{3}{$y^2=2px$，$p>0$、$0\le x\le a$（包括上下两支，$a>0$）。}
\end{minipage}\par\vfill
\noindent\begin{minipage}{\textwidth}\ProblemHeading{习题}{7.15}
计算下列曲线的弧长：
\SubQ{1}{$y=\ln(\cos x)$，$0\le x\le\pi/3$。}
\SubQ{2}{$y=\sqrt{25-x^2}$，$0\le x\le5\sqrt2/2$。}
\end{minipage}\par\vfill
\newpage
\noindent\begin{minipage}{\textwidth}\ProblemHeading{习题}{7.16}
设心脏线的极坐标方程为 $r=a(1+\cos\theta)$，$a>0$。计算：
\SubQ{1}{心脏线的长度。}
\SubQ{2}{其围成的区域绕极轴旋转一周所成的立体的体积。}
\SubQ{3}{该旋转体的侧面积。}
\end{minipage}\par\vfill
\noindent\begin{minipage}{\textwidth}\ProblemHeading{习题}{7.17}
求曲线 $y=\sin x$（$0\le x\le\pi$）绕 $x$ 轴旋转一周所产生的曲面的面积。
\end{minipage}\par\vfill
\noindent\begin{minipage}{\textwidth}\ProblemHeading{习题}{7.18}
设曲线 $y=\sin x$ 与直线 $x=0$、$x=\pi/2$、$y=t$（$0\le t\le1$）所围部分的总面积为 $s(t)$。求 $s(t)$ 的最大值和最小值。
\end{minipage}\par\vfill
\newpage
\noindent\begin{minipage}{\textwidth}\ProblemHeading{习题}{7.19}
已知 $f$ 在 $[0,1]$ 上可导，并满足 $xf^{\prime}(x)=f(x)+3x^2$。曲线 $y=f(x)$ 与直线 $x=0$、$x=1$、$y=0$ 围成的平面图形绕 $x$ 轴旋转一周所成的旋转体的体积达到最小值时，求该平面图形的面积。
\end{minipage}\par\vfill
\noindent\begin{minipage}{\textwidth}\ProblemHeading{例题}{7.20}
一根长为 $l>0$、质量为 $M$ 的均匀细棒 $AB$，在 $A$ 端外侧的延长线上有质量为 $m$ 的质点 $C$，且 $CA=a>0$。求细棒对质点的万有引力，万有引力常数为 $G$。
\end{minipage}\par\vfill
\noindent\begin{minipage}{\textwidth}\ProblemHeading{例题}{7.21}
将质量为 $m$ 的物体从距离地心为 $h$ 的地点移至无穷远处，求克服地球引力所做的功。设地球质量为 $M$，万有引力常数为 $G$，$h$ 不小于地球半径。
\end{minipage}\par\vfill
\newpage
\noindent\begin{minipage}{\textwidth}\ProblemHeading{例题}{7.22}
将物体从地面竖直向空中发射，至少需要多大的初速度才能脱离地球引力？设地球半径为 $R$，地表重力加速度为 $g$，忽略空气阻力与地球自转。
\end{minipage}\par\vfill
\noindent\begin{minipage}{\textwidth}\ProblemHeading{习题}{7.23}
底面半径为 $5\,\mathrm m$、高为 $10\,\mathrm m$ 的圆柱形水桶盛满水。将水全部排至桶口，求所需做的功。水的密度为 $\rho$，重力加速度为 $g$。
\end{minipage}\par\vfill
\noindent\begin{minipage}{\textwidth}\ProblemHeading{习题}{7.24}
一根无限长的均匀细棒，线密度为 $\rho$。距离细棒为 $a>0$ 处有一个单位质量的质点。求细棒对质点的万有引力，万有引力常数为 $G$。
\end{minipage}\par\vfill
\newpage
\noindent\begin{minipage}{\textwidth}\ProblemHeading{习题}{7.25}
直径为 $1\,\mathrm m$、高为 $2\,\mathrm m$ 的直立圆柱桶盛满水，桶底有直径为 $1\,\mathrm{cm}$ 的小孔。水流速度为 $v=0.6\sqrt{2gh}$，其中 $h$ 为瞬时水深，$g$ 为重力加速度。求全部水流完所需的时间。
\end{minipage}\par\vfill
\noindent\begin{minipage}{\textwidth}\ProblemHeading{例题}{7.26（改写）}
设 $\lambda>0$，$m\in\mathbb N^*$，$\xi\ge m$。利用积分估计证明\[\left|\sum_{n=m}^{\lfloor\xi\rfloor}n^\lambda-\frac{\xi^{\lambda+1}-m^{\lambda+1}}{\lambda+1}\right|\le\xi^\lambda.\]并据此给出该和式在固定 $m$、$\xi\to+\infty$ 时的渐近估计。
\end{minipage}\par\vfill
\noindent\begin{minipage}{\textwidth}\ProblemHeading{例题}{7.27（改写）}
设 $T(\xi)=\sum_{1\le n\le\xi}\ln n$（$\xi\ge1$）。证明 $|T(\xi)-\xi\ln\xi+\xi-1|\le\ln\xi$，并推导 $n^{n-1}e^{-n+1}\le n!\le n^{n+1}e^{-n+1}$。
\end{minipage}\par\vfill
\newpage
\noindent\begin{minipage}{\textwidth}\ProblemHeading{例题}{7.28（改写）}
证明极限 $\gamma=\lim_{n\to\infty}(\sum_{k=1}^n1/k-\ln n)$ 存在，并证明 $\sum_{k=1}^n1/k=\ln n+\gamma+O(1/n)$。
\end{minipage}\par\vfill
\noindent\begin{minipage}{\textwidth}\ProblemHeading{例题}{7.29}
求级数 $\sum_{k=1}^{\infty}(-1)^{k-1}/k$ 的和。
\end{minipage}\par\vfill
\noindent\begin{minipage}{\textwidth}\ProblemHeading{例题}{7.30（改写）}
设 $f$ 在 $[1,+\infty)$ 上非负且递减。证明级数 $\sum_{n=1}^{\infty}f(n)$ 与反常积分 $\int_1^{+\infty}f(x)\,dx$ 同敛散。
\end{minipage}\par\vfill
\newpage
\noindent\begin{minipage}{\textwidth}\ProblemHeading{例题}{7.31（改写）}
完成下列不等式问题：
\SubQ{1}{设 $0<a<b$，证明 $\dfrac2{a+b}<\dfrac{\ln b-\ln a}{b-a}<\dfrac12(1/a+1/b)$。}
\SubQ{2}{对 $n\in\mathbb N^{*}$，证明 $\dfrac1{n+1/2}<\ln(1+1/n)<\dfrac12(1/n+1/(n+1))$。}
\end{minipage}\par\vfill
\noindent\begin{minipage}{\textwidth}\ProblemHeading{例题}{7.32}
设 $\varphi$ 在 $[0,+\infty)$ 上连续且严格递增，$\varphi(0)=0$。对 $a>0$、$0<b<\lim_{x\to+\infty}\varphi(x)$，证明 Young 不等式 $ab\le\int_0^a\varphi(x)\,dx+\int_0^b\varphi^{-1}(y)\,dy$，并证明等号当且仅当 $b=\varphi(a)$ 时成立。
\end{minipage}\par\vfill
\noindent\begin{minipage}{\textwidth}\ProblemHeading{例题}{7.33（改写）}
设 $p,q>1$，$1/p+1/q=1$。证明下列不等式，并分别确定等号成立的条件：
\SubQ{1}{对 $a,b>0$，$ab\le a^p/p+b^q/q$。}
\SubQ{2}{对 $A,B>0$，$A^{1/p}B^{1/q}\le A/p+B/q$。}
\end{minipage}\par\vfill
\newpage
\noindent\begin{minipage}{\textwidth}\ProblemHeading{例题}{7.34（改写）}
设 $(a_1,\ldots,a_n)$、$(b_1,\ldots,b_n)$ 为两组各自不全为零的非负数。
\SubQ{1}{$p,q>1$、$1/p+1/q=1$。证明 Hölder 不等式\[\sum_{i=1}^na_ib_i\le\left(\sum_{i=1}^na_i^p\right)^{1/p}\left(\sum_{i=1}^nb_i^q\right)^{1/q},\]并证明等号成立当且仅当存在常数 $\lambda>0$ 使 $a_i^p=\lambda b_i^q$（$i=1,\ldots,n$）。}
\SubQ{2}{证明 Cauchy–Schwarz 不等式 $(\sum_{i=1}^n a_i b_i)^2\le(\sum_{i=1}^n a_i^2)(\sum_{i=1}^n b_i^2)$，并确定等号成立的条件。}
\end{minipage}\par\vfill
\noindent\begin{minipage}{\textwidth}\ProblemHeading{习题}{7.35}
设 $a,b\ge1$。证明 $ab\le e^{a-1}+b\ln b$，并讨论等号成立的条件。
\end{minipage}\par\vfill
\noindent\begin{minipage}{\textwidth}\ProblemHeading{习题}{7.36}
证明对任意 $a>0$，$\sum_{k=1}^{\infty}\dfrac1{k+a}\sqrt{a/k}<\pi$。
\end{minipage}\par\vfill
\newpage
\noindent\begin{minipage}{\textwidth}\ProblemHeading{习题}{7.37（校订）}
设 $f,g\in C[a,b]$，$p,q>1$、$1/p+1/q=1$。证明 Hölder 不等式\[\int_a^b|f(x)g(x)|\,dx\le\left(\int_a^b|f(x)|^p\,dx\right)^{1/p}\left(\int_a^b|g(x)|^q\,dx\right)^{1/q}.\]证明等号当且仅当 $f\equiv0$、$g\equiv0$，或存在 $B>0$ 使 $|f|^p=B|g|^q$ 三种情况之一成立。
\end{minipage}\par\vfill
\noindent\begin{minipage}{\textwidth}\ProblemHeading{习题}{7.38}
设 $f,g\in C[a,b]$，$p\ge1$。证明 Minkowski 不等式\[\left(\int_a^b|f+g|^p\,dx\right)^{1/p}\le\left(\int_a^b|f|^p\,dx\right)^{1/p}+\left(\int_a^b|g|^p\,dx\right)^{1/p}.\]
\end{minipage}\par\vfill
\noindent\begin{minipage}{\textwidth}\ProblemHeading{习题}{7.39}
设 $0<p<1$。证明存在常数 $\beta$，使 $\sum_{k=1}^n1/k^p=n^{1-p}/(1-p)+\beta+O(n^{-p})$。
\end{minipage}\par\vfill
\newpage
\noindent\begin{minipage}{\textwidth}\ProblemHeading{习题}{7.40}
利用面积原理研究和式 $\sum_{k=2}^n\ln\ln k$ 在 $n\to\infty$ 时的渐近行为。
\end{minipage}\par\vfill
\noindent\begin{minipage}{\textwidth}\ProblemHeading{习题}{7.41}
证明存在常数 $\alpha$，使 $\sum_{k=1}^n(\ln k)/k=\frac12\ln^2n+\alpha+O((\ln n)/n)$。
\end{minipage}\par\vfill
\noindent\begin{minipage}{\textwidth}\ProblemHeading{习题}{7.42}
利用 $\sum_{k=1}^n(\ln k)/k=\frac12\ln^2n+\alpha+O((\ln n)/n)$，证明\[\sum_{k=1}^{\infty}(-1)^k\frac{\ln k}{k}=\left(\gamma-\frac12\ln2\right)\ln2,\]其中 $\gamma=\lim_{n\to\infty}(\sum_{k=1}^n1/k-\ln n)$ 为 Euler 常数。
\end{minipage}\par\vfill
\newpage
\noindent\begin{minipage}{\textwidth}\ProblemHeading{习题}{7.43}
设 $\varphi$ 在 $[0,+\infty)$ 上连续且严格递增，$\varphi(0)=0$。对 $a>0$、$0<b<\lim_{x\to+\infty}\varphi(x)$，给出 Young 不等式 $ab\le\int_0^a\varphi(x)\,dx+\int_0^b\varphi^{-1}(y)\,dy$ 的分析证明，不借助面积图。
\end{minipage}\par\vfill
\noindent\begin{minipage}{\textwidth}\ProblemHeading{习题}{7.44}
设 $a_1\ge a_2\ge\cdots\ge a_n\ge0$。证明：
\SubQ{1}{若 $f(0)=0$、$f^{\prime}(0)\ge0$，且 $f^{\prime}$ 在 $[0,a_1]$ 上连续递增，则 $f(\sum_{k=1}^n(-1)^{k+1}a_k)\le\sum_{k=1}^n(-1)^{k+1}f(a_k)$。}
\SubQ{2}{当 $r>1$ 时，$(\sum_{k=1}^n(-1)^{k+1}a_k)^r\le\sum_{k=1}^n(-1)^{k+1}a_k^r$。}
\end{minipage}\par\vfill
\noindent\begin{minipage}{\textwidth}\ProblemHeading{习题}{7.45}
计算 $\lim_{n\to\infty}\frac1n\sum_{k=1}^n(\lfloor2n/k\rfloor-2\lfloor n/k\rfloor)$。
\end{minipage}\par\vfill
\newpage
\noindent\begin{minipage}{\textwidth}\ProblemHeading{习题}{7.46}
设 $-f$ 是 $[a,b]$ 上的凸函数，$f(a)=f(b)=0$，$f^{\prime}(a)=\alpha>0$，$f^{\prime}(b)=\beta<0$。证明 $0\le\int_a^b f(x)\,dx\le\frac12\alpha\beta(b-a)^2/(\beta-\alpha)$。
\end{minipage}\par\vfill
\noindent\begin{minipage}{\textwidth}\ProblemHeading{例题}{7.47}
求极限 $\lim_{n\to\infty}(1+1/n)^{n^2}\dfrac{n!}{n^n\sqrt n}$。
\end{minipage}\par\vfill
\noindent\begin{minipage}{\textwidth}\ProblemHeading{习题}{7.48}
当 $n\to\infty$ 时，估计 $\binom{2n}n$ 的无穷大阶，给出一个等价表达式。
\end{minipage}\par\vfill
\newpage
\noindent\begin{minipage}{\textwidth}\ProblemHeading{习题}{7.49}
证明 $\lim_{n\to\infty}\sqrt n\int_{-1}^1(1-x^2)^n\,dx=\sqrt\pi$。
\end{minipage}\par\vfill
\noindent\begin{minipage}{\textwidth}\ProblemHeading{习题}{7.50}
证明 $\lim_{n\to\infty}(-1)^n\binom{-1/2}n\sqrt n=1/\sqrt\pi$，其中广义二项式系数定义为 $\binom\alpha n=\alpha(\alpha-1)\cdots(\alpha-n+1)/n!$。
\end{minipage}\par\vfill
\noindent\begin{minipage}{\textwidth}\ProblemHeading{习题}{7.51}
设 $\gamma=\lim_{n\to\infty}(\sum_{k=1}^n1/k-\ln n)$ 为 Euler 常数。证明 $\lim_{n\to\infty}\sqrt n\prod_{k=1}^n\dfrac{e^{1-1/k}}{(1+1/k)^k}=\dfrac{\sqrt{2\pi}}{e^{1+\gamma}}$。
\end{minipage}\par\vfill
\newpage
\noindent\begin{minipage}{\textwidth}\ProblemHeading{习题}{7.52}
设实数列 $x_n$ 满足 $x_nx_{n+1}=n$（$n\ge1$），且 $\lim_{n\to\infty}x_n/x_{n+1}=1$。证明 $\pi x_1^2=2$。
\end{minipage}\par\vfill
\noindent\begin{minipage}{\textwidth}\ProblemHeading{习题}{7.53}
利用不等式 $1-x^2\le e^{-x^2}$（$0\le x\le1$）和 $e^{-x^2}\le1/(1+x^2)$（$x\ge0$），证明 $\int_0^{+\infty}e^{-x^2}\,dx=\sqrt\pi/2$。
\end{minipage}\par\vfill
\noindent\begin{minipage}{\textwidth}\ProblemHeading{习题}{7.54}
对 $n\in\mathbb N^*$，令 $S_n=1+\sum_{j=1}^{n-1}\prod_{k=1}^j\dfrac{n-k}{n+k+1}$。证明 $\lim_{n\to\infty}S_n/\sqrt n=\sqrt\pi/2$。
\end{minipage}\par\vfill
